\documentclass[11pt]{article}
\usepackage{amsmath,amssymb,hyperref}
\title{Source-qualified admission of an irreversible handoff}
\author{}
\date{}
\begin{document}
\maketitle
Consider an inspection station with an outstanding old holding duty and a
report due at tick $D$. The station is occupied by a handoff for $h>0$ ticks
starting at $s$. Until completion both the incumbent and candidate holder
are needed. At completion the duty is settled at a secure station, and no
later candidate availability is required. The original holder lasts to $O$.
Holding the duty without transfer is a productive fallback precisely when
$O\ge D$; it then permits the on-time report. A successful nominal transfer
settles the outstanding duty early. A failed attempted transfer is not
abortable or recoverable within this contract.

A source law $(F,L)$ includes no-loss behavior and every loss time $t\ge F$;
the bit observed before the start cut says whether the loss notice, emitted
at $t-L$, has arrived. Notices emitted before the time origin are assumed
recorded. The loss starts at tick $t$; support on $[s,s+h)$ is enough. At the
same tick, notice arrives before dispatch, and completion can coincide with
loss. Missing notices, uncertain sensor delay, interrupted transfers,
unmodeled common causes, and arbitrating among several concurrent starts
are excluded. The law family is a union of trace alternatives, not an
independent product of separately measured floors and leads.

Controlled removal can be prevented by an effector fence until tick $E$.
If controlled removal otherwise becomes possible at $C$, its floor becomes
$\max(C,E)$. Uncommanded failure from $P$ is not delayed by that fence.
When both arms have the same warning law, the combined floor is
$\min(\max(C,E),P)$. When the lead differs, retain the two source pairs. A
snapshot of the epoch at dispatch is not a fence. This needs authority over
the actual controlled effect, including already queued withdrawals. The
model cannot establish that authority for hardware.

\paragraph{Boundary.}
For a finite nonempty family $\mathcal L$ at proposed start $s$, export
\[
 U_s=\min_{(F,L)\in\mathcal L}\max\{F,s+L+1\},
 \qquad W_s=\bigvee_{(F,L)\in\mathcal L}(F\le s+L).
\]
Use the frozen decision rule: on silence initiate; on warning retain the
old holder. Then guaranteed productive admission of that rule is equivalent
to $s+h\le\min(O,D,U_s)$ and $\neg W_s\lor D\le O$.
Below the first silent threshold every loss was already announced; the
threshold itself is realized by some source arm. Its trace has the same
read at $s$ as no loss, preventing a receiver that chooses nominal
dispatch from abandoning only that trace. On the warning branch the old
holder has to cover the report; no elapsed transfer time can be credited.
The interface composes under source union by minimum and disjunction.
Replacing a source while retaining $s,h,O,D$ only needs those tagged
branches. A worsening can invalidate the conclusion; a bound on controlled
withdrawal alone cannot exclude physical failure.

\paragraph{Separations and limits.}
At $s=0,h=4,O=D=5$, source pairs $(1,3)$ and $(5,0)$ admit dispatch on
silence and successful fallback on warning. The lossy pair $(1,0)$ rejects
that design by combining the early onset of one source with the short
warning of the other. At $s=0,h=1,O=1,D=2$, pairs $(1,1)$ and $(4,0)$
violate the frozen warning rule, whereas $(1,0)$ incorrectly approves
that fixed rule. That second example is not an impossibility for all
policies. If no productive fallback exists, the complete two-action
controller must dispatch on every reachable observation and may ignore a
notice when safe. This can succeed exactly when
$s+h\le\min(O,D,\min F)$. If productive fallback exists, the controller
can always retain the holder, but only robust dispatch on the nominal
branch settles early. For that stronger service objective, the warning
correlation in $U_s$ is necessary to recover the example.

In another witness, the candidate is healthy at tick zero, a two-tick
handoff is required, the holder expires at two, and the report is due at
four. Allowing controlled withdrawal at tick one breaks the handoff under
a snapshot-only policy. A true fence through tick two works when physical
loss cannot occur before then; changing only the independent physical
loss floor to one breaks the same guarantee.

\paragraph{Comparison.}
A direct finite temporal policy solver with both source pairs and the
same timing order ties. The benefit found is a sufficient and exact
small receiver signature for this source family, and a counterexample
to two unsafe or incomplete sources of compiled bounds. This is a
restricted temporal controllability calculation; see
\href{https://drops.dagstuhl.de/entities/document/10.4230/LIPIcs.TIME.2017.8}{Cairo and Rizzi on dynamic controllability}
and
\href{https://arxiv.org/abs/1608.04132}{Cairo, Comin, and Rizzi on same-instant reaction}.
The present check does not claim an encoding into their specific network
normal forms or a win against those algorithms. Script totals:
387,072 one-source admission cells, 172,800 two-source judgments under
the frozen rule and 90,000 unrestricted two-action judgments. When
fallback is productive, 3,016 cells are strict false rejections after
independent worst-case projection. The five Lean lemmas check the
definitional arithmetic and not correspondence to physical facilities.
\end{document}
